\documentclass[11pt,a4paper]{article}
\usepackage[margin=25mm]{geometry}
\usepackage[T1]{fontenc}
\usepackage{amsmath,amssymb,amsthm,mathtools,booktabs,microtype,array}
\usepackage[colorlinks=true,linkcolor=blue,citecolor=blue,urlcolor=blue]{hyperref}
\newtheorem{theorem}{Theorem}[section]
\newtheorem{lemma}[theorem]{Lemma}
\newtheorem{proposition}[theorem]{Proposition}
\newtheorem{corollary}[theorem]{Corollary}
\theoremstyle{remark}\newtheorem{remark}[theorem]{Remark}
\newcommand{\R}{\mathbb R}\newcommand{\C}{\mathbb C}
\newcommand{\T}{\mathbb T}\newcommand{\E}{\mathbb E}
\newcommand{\CUT}{\operatorname{CUT}}
\newcommand{\conv}{\operatorname{conv}}
\newcommand{\Tr}{\operatorname{Tr}}
\newcommand{\diag}{\operatorname{diag}}
\newcommand{\diam}{\operatorname{diam}}
\newcommand{\Rea}{\operatorname{Re}}
\title{The arithmetic of sampling static quantum disorder:\\exact aliases, finite witnesses, and clock robustness}
\author{Lluis Eriksson\\\small Independent researcher, Sweden}
\date{15 October 2026 --- working paper}
\begin{document}\maketitle
\begin{abstract}
Whether a dephasing experiment can exclude static Hamiltonian disorder
depends on the arithmetic of its measurement times. We prove that at any
finite set of rationally independent positive times, every real dephasing
semigroup with strictly positive off-diagonal rates is reproduced exactly
by one finite ensemble of static diagonal Hamiltonians. At most
$1+k d(d-1)$ Hamiltonians suffice for $k$ times and $d$ levels; no uniform
energy bound is asserted. We also construct an explicit bounded-energy
static ensemble reproducing the same semigroup on an entire sufficiently
coarse uniform grid. In the opposite direction, every rate matrix outside
the cut cone admits a separating test on finitely many odd harmonics of
one clock period. The proof combines a circle-distance cut representation
with an explicit Fourier-tail bound, and applies to arbitrary static
Hermitian ensembles through their coherence-transfer entries. A qutrit
example uses 20 times and 40,500 shots to bound each error probability by
$5\%$. Consequently, exact finite-time detectability can disappear under
arbitrarily small independent perturbations of the sampling times. A
spectral-width assumption restores robustness with an explicit timing
error budget. All results concern fresh single-use channel experiments,
not interventions within one noise trajectory. The contribution is an
identifiability and experimental-design theory; no quantum environment or
solution of an Erd\H{o}s problem is claimed.
\end{abstract}

\section{Question and relation to existing work}
An ensemble of time-independent Hamiltonians can produce irreversible-
looking coherence decay. A Cauchy-distributed fixed frequency even gives
exact exponential decay. Hamiltonian-ensemble representations and their
Fourier compatibility constraints are established tools
\cite{chen2018,chen2019,chen2021,kropf}. Nevertheless, an incompatibility
of a \emph{whole time-dependent family} need not be visible in a specified
finite collection of its channels. This paper asks which sampling designs
preserve or destroy that distinction.

The answer has three parts. Generic finite menus admit exact static
aliases, carefully related finite menus can exclude every static model,
and coarse infinite grids can admit aliases again. These statements are
compatible: they use different arithmetic relations between times.
Our constructive separation also exposes a discontinuity relevant to
clock calibration. Without an energy bound, independent timing errors
cannot be treated as a uniformly small perturbation of the static null.

The mathematical ingredients have classical antecedents. The torus-density
argument is the continuous Kronecker theorem; the convex-hull and
finite-support arguments are standard finite-dimensional convex analysis
\cite{vdh,rockafellar}. Wrapped Cauchy kernels and circular copulas are
well studied \cite{kato}. Cut-cone and $L^1$ metric representations are
classical \cite{deza}. Our earlier manuscript \cite{parent} gives the
continuous-time cut-cone classification and a diagonal spectral surrogate.
Those results are restated with proofs here, and are not counted as new
independent discoveries. The new results developed here are the exact
finite-menu alias theorem, the explicit simultaneous grid alias, the
complete finite harmonic certificates, and their timing-robustness
consequences. We have not located these combined statements in the cited
literature; this is a scoped priority assessment, not an exhaustive
novelty certificate.

\section{Access model and observables}\label{sec:model}
Fix a distinguished basis of $\C^d$, write $E_{ij}=|i\rangle\langle j|$,
and set $\hbar=1$. A static ensemble has a probability law $\mu$ on
finite Hermitian $d\times d$ matrices and channels
\begin{equation}\label{eq:static}
 \Psi_t(A)=\E_\mu[e^{-itH}A e^{itH}],\qquad t\geq0.
\end{equation}
The law is the same at every queried time; it cannot depend on the fresh
time, input, or readout choice. A shot prepares a fresh system, applies
one such channel, and measures it. The channel model makes no statement
about multiple interventions within the same realization of $H$.

For a real symmetric matrix $\Gamma$ with zero diagonal, the reference is
\begin{equation}\label{eq:semigroup}
 \Phi_t(E_{ij})=e^{-t\Gamma_{ij}}E_{ij}.
\end{equation}
We assume complete positivity: $\Gamma$ is conditionally negative
semidefinite, or equivalently $\Gamma_{ij}=\|v_i-v_j\|^2$ for suitable
real vectors $v_i$ \cite{schoenberg}. In particular, $\Gamma_{ij}\geq0$.
This is a random diagonal unitary channel at each individual time:
if $G$ is a standard real Gaussian vector, the phases
$\sqrt{2t}\langle v_i,G\rangle$ realize \eqref{eq:semigroup}.
This representation changes its phase law with time and does not itself
give a static Hamiltonian law.

For an arbitrary channel define
\begin{equation}\label{eq:T}
 T_{ij}(t)=\langle i|\Psi_t(E_{ij})|j\rangle,\qquad
 C_{ij}(t)=\Rea T_{ij}(t)\quad(i\ne j).
\end{equation}
Population preservation is not assumed for the static null.

\begin{lemma}[Diagonal spectral surrogate]\label{lem:surrogate}
For every law in \eqref{eq:static} there is one joint law of real random
variables $(\varepsilon_1,\ldots,\varepsilon_d)$ such that
\begin{equation}\label{eq:surrogate}
 T_{ij}(t)=\E e^{-it(\varepsilon_i-\varepsilon_j)}
 \quad(i\ne j,\ t\in\R).
\end{equation}
If $\diam\operatorname{spec}(H)\leq B$ almost surely, this law can be
chosen with $|\varepsilon_i-\varepsilon_j|\leq B$ almost surely.
\end{lemma}
\begin{proof}
Conditionally on $H$, independently sample each $\varepsilon_i$ from
the spectral probability measure $\langle i|P_H(\cdot)|i\rangle$.
Then $\E[e^{-it\varepsilon_i}\mid H]=(e^{-itH})_{ii}$, and conditional
independence gives
$\E[e^{-it(\varepsilon_i-\varepsilon_j)}\mid H]
=(e^{-itH})_{ii}\overline{(e^{-itH})_{jj}}$.
The latter is the transfer entry in \eqref{eq:T} before averaging over
$H$. All sampled energies lie in the same conditional spectrum.
\end{proof}

The surrogate reproduces these entries, not necessarily the whole
channel. It is nevertheless enough for every separating statistic below.

\begin{lemma}[Bounded measurement of a transfer entry]\label{lem:phase}
Set $P=E_{ii}+E_{jj}$, $X=E_{ij}+E_{ji}$, and
$Y=-iE_{ij}+iE_{ji}$. Uniformly choose $(A,s)$ from
$\{X,Y\}\times\{-1,1\}$, prepare $(P+sA)/2$, measure $A$ after the
channel, and multiply the outcome by $s$. The resulting variable $O$
belongs to $[-1,1]$ and has mean $C_{ij}(t)$.
\end{lemma}
\begin{proof}
The mean is
$\frac14\Tr[X\Psi_t(X)+Y\Psi_t(Y)]=\Rea T_{ij}(t)$.
The identity uses only linearity and Hermiticity preservation.
The bounded-outcome assertion additionally uses a physical channel,
positive input states, and observables with spectrum in $[-1,1]$.
\end{proof}

\section{Exact aliases on generic finite time menus}\label{sec:alias}
Let $\mathcal C_d=\conv\{zz^*:z\in\C^d,\ |z_i|=1\}$, viewed in
the real affine space of Hermitian matrices with diagonal one. These are
the correlation matrices of random diagonal unitary channels. Its
dimension is $d(d-1)$.

\begin{lemma}[Convex density and the phase body]\label{lem:convex}
If $S$ is dense in a compact set $K\subset\R^m$, then
$\overline{\conv S}=\conv K$ and
$\operatorname{ri}(\conv K)\subseteq\conv S$.
The identity matrix is an interior point of $\mathcal C_d$.
\end{lemma}
\begin{proof}
The closure assertion follows by approximating each vertex of a finite
convex combination. For the relative-interior assertion, take a small
simplex surrounding the point within the affine hull, approximate its
vertices by points of $\conv S$, and retain containment under a small
perturbation. Expanding the vertex combinations gives membership in
$\conv S$.
For the second assertion, Haar-averaging $zz^*$ gives $I$. A supporting
affine functional at $I$ would give a nonnegative trigonometric polynomial
with Haar integral zero, hence the zero polynomial. The distinct
characters $z_i\overline z_j$, $i\ne j$, are linearly independent,
so this functional is constant on the full affine space. Thus there is
no nonconstant supporting functional at $I$.
\end{proof}

\begin{lemma}[Interior of real dephasing data]\label{lem:interior}
If $\Gamma_{ij}>0$ for $i\ne j$, every correlation matrix
$M(t)=(e^{-t\Gamma_{ij}})_{ij}$, $t>0$, is interior to $\mathcal C_d$.
No strict conditional negative definiteness is required.
\end{lemma}
\begin{proof}
Use the Gaussian phase representation from Section~\ref{sec:model}.
Suppose a nonnegative affine functional on $\mathcal C_d$ vanishes at
$M(t)$. Its continuous, nonnegative value on the phase vectors
$z_i(x)=e^{-i\langle v_i,x\rangle}$ has zero integral against a Gaussian
with everywhere-positive density on $\R^r$. It therefore vanishes for
every $x$. Averaging against a Gaussian of any other variance shows that
it vanishes at $M(s)$ for every $s>0$. Since all off-diagonal rates are
positive, $M(s)\to I$ as $s\to\infty$. This contradicts
Lemma~\ref{lem:convex} unless the functional is constant.
\end{proof}

\begin{theorem}[Finite-menu alias theorem]\label{thm:alias}
Let $t_1,\ldots,t_k>0$ satisfy
\begin{equation}\label{eq:independent}
 \sum_{\ell=1}^k n_\ell t_\ell=0,\quad n\in\mathbb Z^k
 \quad\Longrightarrow\quad n=0.
\end{equation}
The closure of the static diagonal-ensemble data at these times is
$\mathcal C_d^k$. Every tuple in the interior of this product is
realized exactly by a static diagonal ensemble supported on at most
$1+k d(d-1)$ Hamiltonians. In particular, this holds for the channels
\eqref{eq:semigroup} whenever all off-diagonal rates are positive.
\end{theorem}
\begin{proof}
The orbit $s\mapsto(e^{-it_1s},\ldots,e^{-it_ks})$ is dense in $\T^k$.
For completeness, the time average of each nonconstant torus character
along this orbit tends to zero by \eqref{eq:independent}. Approximation
by trigonometric polynomials implies convergence of the time averages
of continuous functions to their Haar integrals. A nonempty open set
cannot be missed by the orbit. This is the continuous Kronecker argument.

Choose each energy $\varepsilon_i$ independently as a coordinate in
$\R$, with $\varepsilon_d=0$ as a gauge. The resulting phase arrays
are dense in $(\T^{d-1})^k$. Taking the phase correlation matrix at
each time is a continuous map. Its image closure is the product of
the rank-one phase-correlation sets, whose convex hull is
$\mathcal C_d^k$: independent convex combinations in the factors give
a joint combination by multiplying their weights.
Lemma~\ref{lem:convex} proves exact interior representation by finite
convex combinations of actual static energy vectors. Carath\'eodory's
theorem bounds their number by the affine dimension plus one.
Lemma~\ref{lem:interior} gives the final assertion.
\end{proof}

\begin{corollary}[Statistical indistinguishability]\label{cor:stats}
Under the hypotheses of Theorem~\ref{thm:alias}, no test restricted to
fresh single-use queries at those $k$ times can distinguish the reference
from every static ensemble with type-I and type-II errors whose sum is
less than one. This remains true with arbitrary input states, ancillas,
measurements, and adaptation between queries.
\end{corollary}
\begin{proof}
The alias and reference supply exactly the same channel at every allowed
query. Sequential composition with the experimenter's operations gives
identical outcome laws. Rejection probabilities are therefore equal.
\end{proof}

For a boundary tuple only approximation is asserted. It still excludes
a uniform positive finite-query testing gap: if each entry differs by
at most $\epsilon$, a crude diamond-norm bound for the Schur-channel
difference is $d^2\epsilon$, obtained by summing its elementary
$A\mapsto E_{ii}AE_{jj}$ maps. An $N$-query hybrid argument bounds
outcome total variation by $Nd^2\epsilon/2$.

\begin{remark}
The support bound counts Hamiltonians, not their spectral width.
The construction need not keep energies small as the time menu approaches
a rational relation. This distinction is essential in
Section~\ref{sec:clock}. Zero off-diagonal rates can be handled by first
merging equal vectors $v_i$; the stated theorem avoids that bookkeeping.
\end{remark}

\section{An explicit alias on an infinite coarse grid}\label{sec:grid}
The next construction is independent of the density argument and gives
a bounded-energy law with infinitely many matching observations.
For $0\leq r<1$, the Poisson kernel on the circle is
\begin{equation}\label{eq:poisson}
 P_r(\theta)=\frac{1-r^2}{1-2r\cos\theta+r^2}
 =1+2\sum_{n\geq1}r^n\cos(n\theta).
\end{equation}

\begin{theorem}[Coarse-grid alias]\label{thm:grid}
Let $\Gamma_{ij}>0$ for $i\ne j$, and put
$r_{ij}=e^{-\Delta\Gamma_{ij}}$ for a spacing $\Delta>0$. If
\begin{equation}\label{eq:gridcondition}
 \sum_{i<j}\frac{2r_{ij}}{1+r_{ij}}\leq1,
\end{equation}
then the density, relative to Haar probability measure on $\T^d$,
\begin{equation}\label{eq:gridlaw}
 q(\theta)=1+\sum_{i<j}\bigl(P_{r_{ij}}(\theta_i-\theta_j)-1\bigr)
\end{equation}
defines a static diagonal ensemble reproducing \eqref{eq:semigroup}
at every $t=n\Delta$, $n=0,1,2,\ldots$. Its energies may be chosen in
$[-\pi/\Delta,\pi/\Delta)$, hence with spectral width at most
$2\pi/\Delta$. Writing $M=\binom d2$ and
$\gamma_{\min}=\min_{i<j}\Gamma_{ij}$, the sufficient condition
\begin{equation}\label{eq:coarsebound}
 \Delta\geq\frac{\log(2M-1)}{\gamma_{\min}}
\end{equation}
implies \eqref{eq:gridcondition}; for $d=2$ any positive spacing works.
\end{theorem}
\begin{proof}
Every summand has Haar mean zero. The minimum of $P_r$ is
$(1-r)/(1+r)$, so \eqref{eq:gridcondition} makes $q$ nonnegative.
Distinct positive multiples of distinct roots $e_i-e_j$, with $i<j$,
are distinct Fourier frequencies, even up to sign. Equation
\eqref{eq:poisson} and orthogonality therefore give
\[
 \int e^{-in(\theta_i-\theta_j)}q(\theta)\,dm(\theta)
   =r_{ij}^{|n|}\quad(n\ne0).
\]
Choose centered representatives of the phases and set
$H=\diag(\theta_1,\ldots,\theta_d)/\Delta$. Integer multiples of
$\Delta$ remove the representative ambiguity. All populations are fixed,
so equality of the coherence factors is equality of the whole channels.
Finally, $r_{ij}\leq(2M-1)^{-1}$ implies each summand in
\eqref{eq:gridcondition} is at most $1/M$.
\end{proof}

This sufficient condition is not asserted to be necessary or optimal.
The Poisson kernel itself is standard, as are circular models with
wrapped Cauchy marginals \cite{kato}. The particular sum in
\eqref{eq:gridlaw} is used because its disjoint Fourier supports prescribe
\emph{all} integer moments of \emph{all} pair differences simultaneously.
It need not be a wrapped multivariate Cauchy distribution.

\section{Finite harmonic certificates are complete}\label{sec:finite}
For a subset $S\subseteq\{1,\ldots,d\}$ let
$\delta_S(i,j)=|1_S(i)-1_S(j)|$. The cut cone $\CUT_d$ consists of all
nonnegative combinations of these cut semimetrics. It is a closed,
finitely generated cone. All pair sums below are over $i<j$.

Let $\rho(\theta)=\min_{m\in\mathbb Z}|\theta-2\pi m|\in[0,\pi]$
be geodesic distance to zero on the circle. Its absolutely convergent
Fourier series is
\begin{equation}\label{eq:rho}
 \rho(\theta)=\frac\pi2-\frac4\pi
       \sum_{\substack{n\geq1\\n\ {
m odd}}}\frac{\cos(n\theta)}{n^2}.
\end{equation}

\begin{lemma}[Circular cut representation]\label{lem:circle}
For arbitrary phases $\theta_1,\ldots,\theta_d$, the matrix
$(\rho(\theta_i-\theta_j))_{ij}$ belongs to $\CUT_d$. Hence for every
static ensemble and every $\Delta>0$ the matrix with off-diagonal entries
\begin{equation}\label{eq:D}
 D_{ij}(\Delta)=\frac\pi2-\frac4\pi
  \sum_{n\geq1,\ n\ {
m odd}}\frac{C_{ij}(n\Delta)}{n^2}
\end{equation}
and diagonal zero belongs to $\CUT_d$.
\end{lemma}
\begin{proof}
For the rotating half-circle $I_\alpha=[\alpha,\alpha+\pi)$,
\[
 \rho(\theta_i-\theta_j)=\frac12\int_0^{2\pi}
  |1_{I_\alpha}(\theta_i)-1_{I_\alpha}(\theta_j)|\,d\alpha.
\]
Partitioning the integral by its finitely many induced subsets gives a
nonnegative cut combination. Apply this identity to
$\theta_i=\Delta\varepsilon_i\pmod{2\pi}$ from
Lemma~\ref{lem:surrogate} and take expectations.
Absolute convergence justifies \eqref{eq:D}.
\end{proof}

For the reference semigroup these entries are $F(\Delta\Gamma_{ij})$,
where
\begin{align}
 F(s)&=\frac\pi2-\frac4\pi\sum_{n\geq1,\ n\ {
m odd}}
                          \frac{e^{-ns}}{n^2},\label{eq:F}\\
 F'(s)&=\frac2\pi\log\coth(s/2),\qquad
 F(s)=\frac2\pi s\log(1/s)+O(s)\quad(s\downarrow0).\label{eq:Fasym}
\end{align}
Here $F(0)=0$. To verify the derivative, sum
$\sum_{n\ {
m odd}}e^{-ns}/n=\frac12\log\coth(s/2)$ and integrate
from zero. Consequently, for every fixed $\gamma\geq0$,
\begin{equation}\label{eq:ratelimit}
 \lim_{\Delta\downarrow0}\frac{\pi F(\Delta\gamma)}
                      {2\Delta\log(1/\Delta)}=\gamma.
\end{equation}

\begin{proposition}[Continuous-time classification, with a circle proof]
\label{prop:classification}
The real semigroup \eqref{eq:semigroup} is static exactly when
$\Gamma\in\CUT_d$. Necessity still holds if only all its real
off-diagonal transfer entries agree with an arbitrary static Hermitian
ensemble for every $t\geq0$.
\end{proposition}
\begin{proof}
Necessity follows from Lemma~\ref{lem:circle}, positive rescaling,
\eqref{eq:ratelimit}, and closedness of the cut cone. For sufficiency
write $\Gamma=\sum_S a_S\delta_S$ with $a_S\geq0$. Let independent
centered Cauchy variables $Z_S$ have characteristic functions
$\E e^{-itZ_S}=e^{-a_S|t|}$ and set
$\varepsilon_i=\sum_S1_S(i)Z_S$. Then
$\E e^{-it(\varepsilon_i-\varepsilon_j)}=e^{-|t|\Gamma_{ij}}$,
which gives the required diagonal ensemble.
\end{proof}

For $K\geq1$ define $\mathcal O_K=\{1,3,\ldots,2K-1\}$ and
\begin{equation}\label{eq:tail}
 b_K=\frac4\pi\sum_{n\geq2K+1,\ n\ {\rm odd}}\frac1{n^2}
     =\frac\pi2-\frac4\pi\sum_{n\in\mathcal O_K}\frac1{n^2}.
\end{equation}
Replace the infinite sum in \eqref{eq:D} by $n\in\mathcal O_K$ to obtain
$D^{(K)}_{ij}(\Delta)$. Uniformly over every static energy law,
$|D^{(K)}_{ij}-D_{ij}|\leq b_K$.

\begin{theorem}[Complete finite harmonic separation]\label{thm:finite}
Let $\Gamma\notin\CUT_d$. Choose a real pair functional $w$ such that
\begin{equation}\label{eq:separator}
 \sum_{i<j}w_{ij}\delta_S(i,j)\leq0\quad\hbox{for every }S,
 \qquad \sum_{i<j}w_{ij}\Gamma_{ij}>0,
\end{equation}
and put $W=\sum_{i<j}|w_{ij}|$. There exist $\Delta>0$ and a finite
$K$ such that the statistic
\begin{equation}\label{eq:S}
 S_{K,\Delta}=\sum_{i<j}w_{ij}D^{(K)}_{ij}(\Delta)
\end{equation}
has null upper bound $U=b_KW$ for every static Hermitian ensemble,
while its reference value $V$ satisfies $g:=V-U>0$.
Only the $K$ times $\Delta,3\Delta,\ldots,(2K-1)\Delta$ are needed.
Conversely, a rate matrix in $\CUT_d$ cannot be excluded by any valid
static-null test on the whole reference family.
\end{theorem}
\begin{proof}
A separator exists because $\CUT_d$ is closed and convex.
The null bound follows from Lemma~\ref{lem:circle} and
$|w\cdot(D^{(K)}-D)|\leq b_KW$. By \eqref{eq:ratelimit},
$w\cdot(F(\Delta\Gamma_{ij}))$ is positive for sufficiently small
$\Delta$. Fix such a spacing and increase $K$ until twice the truncation
bound $b_KW$ is smaller than this positive number. Then $V>b_KW$.
The converse follows from Proposition~\ref{prop:classification}.
\end{proof}

This proves existence and supplies computable certificates; it does not
claim polynomial complexity in dimension or an optimal number of times.
For any proposed $w$, its validity can be checked on the finitely many
cuts. Thus the general certificate does not require trusting the output
of a cone-membership solver.

\section{Finite-shot experiment and a certified example}\label{sec:shots}
Write the statistic in the form
\begin{equation}\label{eq:amplitude}
\begin{aligned}
 S_{K,\Delta}&=c+\sum_{i<j,\ n\in\mathcal O_K}a_{ijn}C_{ij}(n\Delta),
 \qquad c=\frac\pi2\sum_{i<j}w_{ij},\\
 a_{ijn}&=-\frac{4w_{ij}}{\pi n^2},\qquad
 A=\sum|a_{ijn}|=\frac{4W}{\pi}\sum_{n\in\mathcal O_K}\frac1{n^2}.
\end{aligned}
\end{equation}
Choose $(i,j,n)$ with probability $|a_{ijn}|/A$ and perform the
four-phase experiment of Lemma~\ref{lem:phase} at $n\Delta$.
Set $Z=\operatorname{sgn}(a_{ijn})O$. Then $|Z|\leq1$ and
$\E Z=(S_{K,\Delta}-c)/A$.

\begin{proposition}[Shot guarantee]\label{prop:shots}
Given $g=V-U>0$, reject the static null when
$c+A\overline Z>(U+V)/2$. For independent fresh shots, each error
probability is at most
\begin{equation}\label{eq:shots}
 \exp\!\left(-\frac{Ng^2}{8A^2}\right).
\end{equation}
The null guarantee also holds with predictable drift if, conditionally
on the past, the current Hamiltonian law is chosen before the independent
random measurement settings and obeys the same static null bound.
\end{proposition}
\begin{proof}
The normalized separation is $g/A$ and each decision threshold is at
least $g/(2A)$ from the corresponding mean. Hoeffding's inequality for
$[-1,1]$-valued variables gives \eqref{eq:shots}. Its conditional
exponential-moment proof gives the stated drift version.
\end{proof}

Consider the qutrit rates
\begin{equation}\label{eq:qutrit}
 \Gamma=\begin{pmatrix}0&1&4\\1&0&1\\4&1&0\end{pmatrix}.
\end{equation}
They arise from the generator $-[D,[D,\cdot]]$ with
$D=\diag(0,1,2)$ and so define physical channels. The triangle
functional $w_{13}=1$, $w_{12}=w_{23}=-1$ is nonpositive on every cut,
but $w\cdot\Gamma=2$. Thus the whole family has no static representation.

\begin{corollary}[Twenty-time certificate]\label{cor:20}
For \eqref{eq:qutrit}, take $K=20$ and $\Delta=13/100$ in the time
units fixed by \eqref{eq:qutrit}. Then
\begin{align*}
 b_{20}&\in(0.01591218,0.01591219),&
 A&\in(4.66465243,4.66465244),\\
 g&\in(0.11365,0.11366).&&
\end{align*}
The maximum queried time is $39\Delta=5.07$. A total of 40,500 shots
is sufficient for each error probability to be less than $0.05$.
\end{corollary}
\begin{proof}
Here $W=3$, $c=-\pi/2$, and the exact margin is
\begin{equation}\label{eq:qmargin}
 g=\frac4\pi\sum_{n\in\mathcal O_{20}}
       \frac{2e^{-13n/100}-e^{-52n/100}+3}{n^2}-2\pi.
\end{equation}
The companion verifier brackets this finite expression using rational
arithmetic: alternating arctangent sums in Machin's formula bracket
$\pi$, and positive Taylor sums with a geometric tail bracket each
exponential. The displayed conservative bounds imply
$40500\,g^2/(8A^2)>3$; the positive series for $e^3$ gives $e^3>20$.
Equation~\eqref{eq:shots} completes the proof.
\end{proof}

The same example has an exact alias on \emph{every} grid time
$n\log4$. Indeed $r_{12}=r_{23}=1/4$ and $r_{13}=1/256$, and
\eqref{eq:gridcondition} becomes
$4/5+2/257=1038/1285<1$.
The density \eqref{eq:gridlaw} is therefore at least $247/1285$.
This gives an explicit pair of physically different models whose entire
coarse-grid channel data are identical, although the 20-time certificate
excludes every static model on its finer, deliberately related menu.

\section{Clock fragility and a quantitative repair}\label{sec:clock}
The set of time vectors satisfying \eqref{eq:independent} has full
Lebesgue measure and is dense: its complement is a countable union of
integer hyperplanes.

\begin{corollary}[No unrestricted independent-jitter guarantee]
\label{cor:jitter}
Fix a real dephasing semigroup with positive off-diagonal rates and any
finite nominal menu of positive times. Every open neighborhood of that
menu contains actual timings at which one static finite ensemble
reproduces the reference channels exactly. Consequently, no test has a
uniform positive separation over an unrestricted static null and an
independent nonzero uncertainty interval for each time setting.
\end{corollary}
\begin{proof}
Choose a rationally independent vector in the uncertainty box and apply
Theorem~\ref{thm:alias} and Corollary~\ref{cor:stats}.
\end{proof}

The corollary concerns fixed but independently uncertain menu entries.
It does not assert that all clock errors destroy detection. A common
multiplicative clock error preserves harmonic relations. Nor does it
cover a continuum of freshly randomized times whose law is independent
of the hidden Hamiltonian; that is a different experiment, as in
\cite{parent,qutrit}.

\begin{theorem}[Timing stability]\label{thm:clock}
Suppose the static null satisfies
$\diam\operatorname{spec}(H)\leq B$ almost surely. If every implemented
time differs from its intended time by at most $\eta$, the statistic
\eqref{eq:amplitude} has null upper bound $U+AB\eta$.
Set $G=\max_{i<j}\Gamma_{ij}$. The value under the reference
\eqref{eq:semigroup} is at least $V-AG\eta$.
Thus a positive certified gap remains whenever
\begin{equation}\label{eq:clockgap}
 g_\eta=g-A(B+G)\eta>0.
\end{equation}
Using the midpoint of the adjusted bounds in the decision rule gives
each error probability at most $\exp[-Ng_\eta^2/(8A^2)]$.
\end{theorem}
\begin{proof}
Lemma~\ref{lem:surrogate} and
$|e^{ix}-e^{iy}|\leq|x-y|$ imply
$|C_{ij}(t)-C_{ij}(s)|\leq B|t-s|$.
The reference factors satisfy
$|e^{-t\Gamma_{ij}}-e^{-s\Gamma_{ij}}|\leq G|t-s|$ for nonnegative
times. Sum the biases against the absolute coefficients in
\eqref{eq:amplitude} and apply Proposition~\ref{prop:shots}.
The same argument permits random timing errors bounded almost surely,
provided the setting-independent Hamiltonian assumption is retained.
\end{proof}

For the qutrit example, $G=4$. The sufficient calibration requirement is
$\eta<g/[A(B+4)]$; setting $B=10$ gives
$\eta<0.00174$ in the chosen time units. This is a bound conditional on
a separately justified spectral-width prior, not a bound inferred from
the same finite data. The high-energy aliases in
Theorem~\ref{thm:alias} explain why such an assumption matters.

\section{Verification, interpretation, and open questions}
The companion code supplies an exact rational certificate for
Corollary~\ref{cor:20}, checks the Fourier moments and positivity of the
explicit qutrit grid density, and tests the phase estimator and spectral
surrogate against direct unitary evolution in several dimensions. A
finite linear-programming example illustrates exact interpolation on
two rationally independent times. That numerical interpolation is an
illustration; Theorem~\ref{thm:alias} supplies its existence proof.
No proof-assistant formalization or laboratory validation is claimed.

The results are stronger than simply observing ordinary frequency
aliasing. They concern one positive joint energy law matching many
transition channels simultaneously, and identify an arithmetic contrast
between entire classes of experimental menus. At the same time, the
proof ingredients are elementary and the scope is specific. The null
assumes a fresh single-use channel on each shot. Multi-time interventions,
echo controls within a single trajectory, or measurements of the hidden
Hamiltonian may provide information unavailable in this model. Likewise,
failure of a static explanation can be caused by classical time-dependent
noise; it does not certify nonclassical environmental correlations.

Several quantitative problems remain. The coarse-grid positivity
criterion is only sufficient. The smallest bandwidth of a finite-menu
alias, especially near a harmonic menu, is not determined. Neither the
20-time design nor its shot bound is claimed optimal. Quantitative
Diophantine approximation may bound alias energies, while sharper
trigonometric majorants and statistical estimators may reduce certificate
cost. These are open extensions of the stated theorems, not results
silently assumed in them.

\paragraph{Reproducibility and assistance.}
The release package contains the source, verification code and output,
public source inventory, and a scoped literature audit. The author used
OpenAI Codex for research assistance, mathematical drafting, code, and
verification. AI assistance and any subsequent model assessments do not
constitute human peer review. Publication as a working paper is distinct
from editorial acceptance or validation of a claimed priority.

\begin{thebibliography}{99}
\bibitem{chen2018} H.-B. Chen, C. Gneiting, P.-Y. Lo, Y.-N. Chen,
and F. Nori, Simulating open quantum systems with Hamiltonian ensembles
and the nonclassicality of the dynamics, \emph{Phys. Rev. Lett.} 120,
030403 (2018). \href{https://arxiv.org/abs/1703.09428}{arXiv:1703.09428}.
\bibitem{chen2019} H.-B. Chen, P.-Y. Lo, C. Gneiting, J. Bae,
Y.-N. Chen, and F. Nori, Quantifying the nonclassicality of pure dephasing,
\emph{Nature Communications} 10, 3794 (2019).
\href{https://doi.org/10.1038/s41467-019-11502-4}{doi:10.1038/s41467-019-11502-4}.
\bibitem{chen2021} H.-B. Chen and Y.-N. Chen, Canonical Hamiltonian ensemble
representation of dephasing dynamics and the impact of thermal
fluctuations on quantum-to-classical transition, \emph{Scientific Reports}
11, 10046 (2021).
\href{https://doi.org/10.1038/s41598-021-89400-3}{doi:10.1038/s41598-021-89400-3}.
\bibitem{kropf} C. M. Kropf, Protecting quantum coherences from static
noise and disorder, \emph{Phys. Rev. Research} 2, 033311 (2020).
\href{https://doi.org/10.1103/PhysRevResearch.2.033311}{doi:10.1103/PhysRevResearch.2.033311}.
\bibitem{vdh} J. van der Hoeven, On the computation of limsups,
corrected author manuscript dated 26 February 1996, Section 2.
\url{https://www.texmacs.org/joris/limsup/limsup-corrected.html}.
\bibitem{rockafellar} R. T. Rockafellar, \emph{Convex Analysis},
Princeton University Press (1970), Sections 6 and 17.
\bibitem{kato} S. Kato, C. Ley, S. Loizidou, and K. V. Mardia,
A versatile trivariate wrapped Cauchy copula with applications to
toroidal and cylindrical data, arXiv:2401.10824, version 2 (2024).
\url{https://arxiv.org/abs/2401.10824}.
\bibitem{deza} M. Deza and M. Laurent, \emph{Geometry of Cuts and
Metrics}, Springer (1997).
\href{https://doi.org/10.1007/978-3-642-04295-9}{doi:10.1007/978-3-642-04295-9}.
\bibitem{schoenberg} I. J. Schoenberg, Metric spaces and positive
definite functions, \emph{Trans. Amer. Math. Soc.} 44, 522--536 (1938).
\href{https://doi.org/10.1090/S0002-9947-1938-1501980-0}{doi:10.1090/S0002-9947-1938-1501980-0}.
\bibitem{parent} L. Eriksson, Static Hamiltonian disorder: exact
classification, global obstructions, and phase-cycled witnesses,
unpublished manuscript, 14 October 2026. Private AIRR deposit
SUB-2CDB76AE0F389DD5, revision 1. The deposited artifact is distinct
from the present manuscript.
\bibitem{qutrit} L. Eriksson, A sharp two-rate reset witness against static qutrit
disorder, AIRR working paper ARR-2026-131G33B01486YA04, version 1 (2026).
\url{https://airr.science/papers/ARR-2026-131G33B01486YA04/versions/v1/}.
\end{thebibliography}
\end{document}
